\section{方法比较与灵敏度分析}
\label{sec:compare}

\subsection{方法对照汇总}

本文对每个问题都实现了多种方法并做对照（\cref{tab:allcmp}）。
需要强调：\textbf{所有方法都在同一“因果口径”下评价}——0:00 只用历史信息（与 0:00 已发布的预报）
制定计划，日内按第 \ref{sec:q2} 节的实时平衡控制规则结算（式 \eqref{eq:clip}--\eqref{eq:socrec}），
因此表中每一行的数字都是\textbf{可实行的策略}在同一信息结构下的成绩，不存在“事后调整”带来的虚优。

\begin{table}[htbp]
\centering
\caption{全部方法的对照汇总（单位：元；问题二起统一采用因果口径实时平衡结算）}
\label{tab:allcmp}
\scriptsize
\begin{tabularx}{\textwidth}{@{}p{1.5cm}Xrrr@{}}
\toprule
问题 & 方法 & 费用 & 相对采用方案 & 结论 \\
\midrule
\multirow{3}{*}{一}
& \textbf{线性规划（采用）} & \textbf{35\,126.95} & --- & 精确最优 \\
& 贪心配对启发式 & 42\,187.93 & $+20.1\text{\%}$ & 忽略时序约束仍更差 \\
& 无储能基线 & 48\,052.05 & $+36.8\text{\%}$ & 储能价值 26.9\% \\
\midrule
\multirow{9}{*}{二}
& \textbf{升级预测模型（采用）} & \textbf{13\,655\,835} & --- & 距下界 136.9 万 \\
& 原依据（同类日均值 $+$ 近 4 日光伏） & 13\,848\,183 & $+1.41\text{\%}$ & 预测依据升级前 \\
& 附件1 典型日曲线计划 & 18\,318\,686 & $+34.1\text{\%}$ & 忽略日际变化 \\
& 无储能 & 18\,155\,061 & $+32.9\text{\%}$ & 储能价值 449.9 万 \\
& 昨日外推 + 报童 $q=0.6$ & 21\,491\,755 & $+57.4\text{\%}$ & 仅用昨日信息 \\
& 昨日外推 + 报童 $q=0.8$ & 21\,789\,429 & $+59.6\text{\%}$ & 裕度过大 \\
& 昨日外推（无裕度） & 22\,515\,792 & $+64.9\text{\%}$ & 朴素外推 \\
& 两阶段随机规划 SAA & 39\,954\,010 & $+192.6\text{\%}$ & 见下文失效分析 \\
& 天眼链式下界 & 12\,286\,531 & $-10.0\text{\%}$ & 不可达 \\
\midrule
\multirow{5}{*}{三}
& \textbf{升级预测、12:00 单次调整（采用）} & \textbf{13\,576\,891} & --- & 距下界 129.0 万 \\
& 全滚动（6/12/18） & 13\,583\,600 & $+0.05\text{\%}$ & 多调无益 \\
& 仅 0:00 计划（不调整） & 13\,655\,835 & $+0.58\text{\%}$ & 调整价值 7.89 万 \\
& 仅 6:00 调整 & 13\,839\,771 & $+1.94\text{\%}$ & 该时刻预报无新信息 \\
& 天眼链式下界 & 12\,286\,531 & $-9.5\text{\%}$ & 不可达 \\
\midrule
\multirow{5}{*}{四}
& \textbf{Q4-2 升级预测 + 平衡规则（采用）} & \textbf{14\,339\,252} & --- & 储能价值 450.6 万 \\
& Q4-2 计划改用预测电价 & 14\,311\,077 & $-0.20\text{\%}$ & 电价未知的代价 $\approx0$ \\
& Q4-2 无储能 & 18\,844\,994 & $+31.4\text{\%}$ & 储能价值对照 \\
& Q4-3 12:00 调整 & 14\,262\,281 & $-0.54\text{\%}$ & 调整价值 7.70 万 \\
& Q4-3 仅 6:00 调整 & 14\,526\,380 & $+1.30\text{\%}$ & 过早动作更差 \\
\bottomrule
\end{tabularx}
\end{table}

由\cref{tab:allcmp} 可见：\textbf{“升级预测模型（分类组合负荷 $+$ 岭回归堆叠光伏）$+$
报童安全裕度 $+$ LP 计划 $+$ 0:00 定下的平衡响应规则（问题三再叠加 12:00 调整）”的组合
在所有可比方案中费用最低}，且距不可达的天眼下界仅 9.5\%。三个值得强调的结论：

\begin{enumerate}[label=(\arabic*), leftmargin=2.2em, itemsep=1pt]
\item \textbf{用统计/机器学习方法提升预测精度，比引入更多预报时刻更有效}：
把 0:00 依据从“近 4 个同类日均值 + 近 4 日光伏均值”升级为“水平$\times$形状分类组合 +
岭回归堆叠”，问题二费用由 13\,848\,183 元降到 \textbf{13\,655\,835 元}（$-1.39\text{\%}$）；
而在此基础上，6:00/18:00 的预报都不值得引入（\cref{tab:q3cmp}）。
\item \textbf{“只在 0:00 决策一次”与“用储能实时兜底”并不矛盾}：平衡响应规则本身在 0:00
就已冻结（式 \eqref{eq:clip}--\eqref{eq:socrec} 不含任何未来信息，也不需要日内重新优化），
它只是把“计划—实际”的偏差按“先调计划充放、再动用储能”的优先级消化。
若改为\textbf{僵硬执行计划}（缺口全部按 5 倍紧急购电），全年费用升至 14\,557\,730 元，
即这条 0:00 定下的规则价值 \textbf{90.2 万元（$+6.60\text{\%}$）}。
\item \textbf{“已知实时电价”并非最优假设}：用 MAE $=0.0367$ 元/kWh 的电价预测替代“已知电价”，
结算按实际电价，全年反而低 0.20\%（\cref{tab:q4cmp}），说明计划阶段过度追价会损害兜底能力。
\end{enumerate}

\textbf{SAA 失效机理分析}。本文将 0:00 计划建模为两阶段随机规划\cite{vtpp}：
以最近 20 天的实测（负载、光伏）曲线为等概率情景，决策计划购电量 $P^0$
使“计划费 $+$ 期望紧急费”最小。结果显示（\cref{tab:allcmp}）该方案全年费用高达
39\,954\,010 元（比采用方案高 192.6\%），其中计划购电费仅 875.3 万元（比采用方案低 32\%）、
紧急购电费却达 3\,120.1 万元（紧急电量 6.93 GWh，是可执行方案的 54 倍）。
原因是\textbf{情景内的储能自由度被重复利用}：在 20 个情景中储能都能从当日初始电量出发、
以各情景自身最优的方式兜底，故情景内紧急购电几乎为零，优化器据此大幅压低 $P^0$；
而真实日的偏差一旦同时差于多数情景，储能已被当日调度消耗，只能购买 5 倍紧急电。
这提示：\textbf{小样本经验情景下的随机规划会产生系统性乐观偏差，
在“计划—紧急”强惩罚机制下反而不如稳健的报童裕度}；
若要用随机规划，必须对“储能可复用性”与“情景外偏差”做显式修正。

\subsection{参数灵敏度分析}

\begin{table}[htbp]
\centering
\caption{灵敏度分析（全部在升级预测口径下重算；单位：元）}
\label{tab:sens}
\scriptsize
\begin{tabular}{@{}p{2.1cm}p{3.4cm}rrr@{}}
\toprule
扰动项 & 变动方式 & 基准 & 变动后 & 相对变化 \\
\midrule
储能放电下限 & 物理下限 1200 $\to$ 2200 / 4000 kWh（1 月） & 1\,237\,127 & 1\,392\,750 / 1\,843\,118 & $+12.6$ / $+49.0\text{\%}$ \\
报童分位偏低 & $q=0.7\to0.5/0.6$ & 13\,655\,835（全年） & 14.79 / 14.05 M & $+8.27$ / $+2.91\text{\%}$ \\
报童分位偏高 & $q=0.7\to0.8/0.9$ & 13\,655\,835 & 13.62 / 13.93 M & $-0.28$ / $+1.98\text{\%}$ \\
充放电效率口径 & 单程 0.9/0.9（往返 81\%）$\to$ 往返 0.9（单程 0.9487） & 13\,655\,835 & 13\,191\,041 & $-3.40\text{\%}$ \\
结算方式 & 实时平衡控制 $\to$ 僵硬执行计划 & 13\,655\,835 & 14\,557\,730 & $+6.60\text{\%}$ \\
偏差计价口径 & 差额计价（B，采用）$\to$ 计划全额 $+$ 罚（A） & 13\,576\,891（问题三） & 13\,641\,510 & $+0.48\text{\%}$ \\
电价信息 & 已知实时电价 $\to$ 用预测电价（结算按实际） & 14\,339\,252（Q4-2） & 14\,311\,077 & $-0.20\text{\%}$ \\
日末残值系数 & $v=0.9\bar p$ $\to$ $v=0$（日末清空） & 1\,218\,659（2 月） & 1\,214\,270 & $-0.36\text{\%}$ \\
日末储电量 & 自由 $\to$ 日末 SOC $=$ 日初（等式） & 1\,218\,659（2 月） & 1\,215\,187 & $-0.28\text{\%}$ \\
安全裕度形式 & 经验分位 $\to$ 正态/分噪声/60 天/90 天窗 & 13\,655\,835 & 13.73 / 13.59 / 13.66 / 13.68 M & $+0.58$ / $-0.52$ / $+0.03$ / $+0.16\text{\%}$ \\
\bottomrule
\end{tabular}
\end{table}

\textbf{灵敏度结论}：
\begin{enumerate}[label=(\arabic*), leftmargin=2.4em, itemsep=1pt]
\item 结论对\textbf{储能放电下限}（$+12.6\text{\%}$ 以上）、\textbf{报童分位偏低}
（$q=0.5$ 时 $+8.27\text{\%}$）、\textbf{结算方式}（$+6.60\text{\%}$）与
\textbf{效率口径}（$-3.40\text{\%}$）较敏感。前两者是本文的决策变量，已用 1 月数据标定；
结算方式与计价口径已在第 \ref{sec:assume} 节假设中显式声明；
效率口径的敏感性说明\textbf{“储能损耗越低、收益越高”}，本文采用的“单程 0.9”
是附件1 给出的最保守读法，故结果偏保守。
\item 结论对\textbf{残值系数}、\textbf{日末储电量约束}、\textbf{偏差计价口径}、
\textbf{电价信息}与\textbf{安全裕度形式}不敏感（均 $<1\text{\%}$），
说明模型不依赖人为参数微调；其中分噪声叠加形式略优 0.52\%，
但它需要分别估计负荷与光伏两套误差分布、参数更多，本文保留理论更直接的净残差经验分位。
\item 全部“复杂化”尝试（典型日曲线、昨日外推、显式季节项、条件分位回归缓冲、
两阶段随机规划）都\textbf{未能降低成本}，说明本文方案已处于该信息结构下的效率前沿。
\end{enumerate}
